\chapter{物理实验、模型拟合与模拟数据}

\section{本章阅读方式}

所有学生快读模型拟合共同框架；物理实验和模拟方向学生精做。目标不是找到最低误差曲线，而是说明模型、参数、残差、不确定度和解释边界，并理解它如何成为机器学习 baseline。

最低交付是一个 Evidence Record、拟合结果和残差/不确定度检查。本章进入 Part III 的作用，是把物理模型变成可解释 baseline，并让参数、残差和误差指标成为后续模型评估的参照。

\section{本章导读：拟合曲线不是终点，解释参数才是目标}

前面几章主要从不同数据类型出发学习观测记录如何变成测量结果。本章再向前走一步：当我们已经有一个候选机制模型时，数据分析就不只是描述曲线形状，而是要用数据约束模型参数，并判断这些参数是否有合理的单位、数量级和物理解释。

这也是物理实验和许多普通预测任务的根本差别。一个多项式或灵活模型可能把曲线贴得很好，但它未必告诉你系统时间常数是多少、残差是否暴露了模型缺陷、误差条是否被正确使用、参数不确定度是否真实反映了观测限制。对科学训练来说，最低误差并不自动等于最好解释。

本章用指数衰减作为最小例子，是因为它能把前向模型、参数搜索、残差诊断、\(\chi^2\)、Monte Carlo 不确定度和 baseline 比较放在同一条线里。学完本章后，你应该能把“我拟合出了一条曲线”升级成“我用某个机制模型得到了一组带边界的参数估计”。

\section{学习目标}

学完本章后，你应该能够：

\begin{itemize}
    \item 理解物理实验数据与纯表格分类数据的差别在于它通常伴随明确的机制模型；
    \item 说清楚模型、参数、残差和 $\chi^2$ 在拟合中的角色；
    \item 用一个简单的指数衰减例子理解“实验观测 $\leftrightarrow$ 前向模型 $\leftrightarrow$ 参数估计”的链条；
    \item 理解 Monte Carlo 重采样为什么能帮助我们估计参数不确定度；
    \item 比较物理模型拟合和简单数据驱动近似各自的优势与局限。
\end{itemize}

进入本章前，建议你已经能写小函数、读懂循环，并理解随机种子为何影响结果，因为后面的模拟和参数扫描都要反复调用这些最小能力。

\section{训练目标与贯穿微型项目}

本章的微型项目是对一个教学版 RC 放电数据集做指数模型拟合。正文负责解释机制模型、参数、残差、\(\chi^2\) 和 Monte Carlo 不确定度；训练材料则要求你交付一份能被复查的拟合记录，而不是只报告一条曲线或一个最佳参数。

\begin{examplebox}[最小物理拟合项目结构]
\begin{lstlisting}[style=terminal]
physics_fit_demo/
  data/
    physics_model_fit_demo.csv
  notebooks/
    fit_decay_model.ipynb
  scripts/
    fit_decay_tau.py
  figures/
    decay_fit_residuals.png
  results/
    model_fit_evidence_record.txt
\end{lstlisting}
\end{examplebox}

项目交付应包含：变量单位、候选模型、参数搜索范围、最佳参数、残差诊断、Monte Carlo 不确定度、baseline 比较和模型假设限制。这个结构会为 Part III 的模型评估和后续项目报告打底。

\section{为什么物理数据分析不能只停留在“相关性”}

在前面的章节中，我们已经见过星表、图像、光谱和时间序列。到了物理实验或模拟数据这一章，分析视角会再向前迈一步：我们不再满足于“看见规律”，而要进一步问“这个规律是否来自某个机制模型”。

例如，一个随时间单调下降的信号，可能来自：

\begin{itemize}
    \item RC 电路放电；
    \item 放射性衰变计数；
    \item 某种冷却过程；
    \item 带有损耗的实验系统。
\end{itemize}

如果我们只用一个灵活的多项式把它拟合得很好，的确可以描述曲线形状，却不一定能回答“系统时间常数是多少”“是否符合理论预期”“参数有没有物理单位”这些更重要的问题。

\begin{defbox}[物理模型拟合]
物理模型拟合的目标，不只是让曲线贴合数据，更是让拟合参数与某个机制、守恒关系或动力学方程建立对应。
\end{defbox}

\section{实验数据的最小结构}

一个典型的实验型数据集，往往可以写成：

\[
\{(x_i, y_i, \sigma_i)\}_{i=1}^N,
\]

其中 $x_i$ 可以是时间、位置、电压、频率或控制变量，$y_i$ 是测量输出，$\sigma_i$ 则表示测量不确定度。

与“只拿特征做预测”的数据思路相比，这里更关键的是：我们通常已经有一个候选函数族

\[
y_{\mathrm{model}}(x;\theta),
\]

并希望通过数据去约束参数 $\theta$。

\paragraph{从“描述曲线”到“解释系统”}

这正是实验型数据和许多普通机器学习表格任务的最大差别之一。对实验数据来说，我们通常不仅想知道：

\begin{itemize}
    \item 曲线看起来长什么样；
    \item 哪个拟合器误差更小；
\end{itemize}

还想进一步知道：

\begin{itemize}
    \item 参数是否对应某个真实物理量；
    \item 该参数的单位和数量级是否合理；
    \item 如果参数变化，系统响应是否仍符合理论关系。
\end{itemize}

也就是说，实验拟合的真正目标常常不是“压低误差”，而是“把数据重新接回模型和机制”。

\section{指数衰减：一个最小物理例子}

本章采用一个教学版 RC 放电数据集。若电容放电满足理想关系，则电压随时间变化为

\[
V(t) = V_0 e^{-t/\tau},
\]

其中：

\begin{itemize}
    \item $V_0$ 是初始电压；
    \item $\tau = RC$ 是系统时间常数；
    \item $t$ 是时间。
\end{itemize}

这条公式之所以适合作为教学起点，是因为它同时具备：

\begin{itemize}
    \item 明确的物理参数；
    \item 清楚的函数形状；
    \item 足够简单的拟合方式；
    \item 可以自然连接到残差分析与 Monte Carlo 不确定度估计。
\end{itemize}

\section{从微分方程到指数解}

RC 放电的核心动力学关系可以写成

\[
\frac{dV}{dt} = -\frac{1}{RC}V.
\]

若令 $\tau = RC$，则有

\[
\frac{dV}{dt} = -\frac{1}{\tau}V.
\]

分离变量并积分：

\[
\frac{dV}{V} = -\frac{dt}{\tau}
\quad \Longrightarrow \quad
\int \frac{dV}{V} = -\int \frac{dt}{\tau},
\]

得到

\[
\ln V = -\frac{t}{\tau} + C.
\]

两边取指数后可写成

\[
V(t) = V_0 e^{-t/\tau}.
\]

这一步非常教材化，因为它说明了：\textbf{实验拟合公式并不是凭空写下来的经验曲线，而是可以从动力学方程推导出来。}

\section{残差与卡方统计：判断“贴合得好不好”的最低标准}

如果观测值为 $y_i$，模型预测为 $y_{\mathrm{model}}(x_i;\theta)$，那么残差定义为

\[
r_i = y_i - y_{\mathrm{model}}(x_i;\theta).
\]

残差不是拟合后的“边角料”，它本身就是数据诊断的核心对象。理想情况下，残差应当：

\begin{itemize}
    \item 没有明显系统偏向；
    \item 大小与误差量级相当；
    \item 不呈现显著的结构性趋势。
\end{itemize}

若每个点带有不确定度 $\sigma_i$，常用的评价量是

\[
\chi^2(\theta) = \sum_{i=1}^{N}
\frac{\left[y_i - y_{\mathrm{model}}(x_i;\theta)\right]^2}{\sigma_i^2}.
\]

在最小二乘或最大似然的高斯误差近似下，寻找较小的 $\chi^2$ 就对应寻找更合理的参数。

\paragraph{为什么还要看缩并卡方}

单独报告 $\chi^2$ 往往还不够，因为样本点数越多，$\chi^2$ 通常也会自然增大。为了更方便比较不同拟合设置，我们常进一步看\textbf{缩并卡方}

\[
\chi^2_\nu = \frac{\chi^2}{\nu},
\]

其中 $\nu = N - p$ 是自由度，$N$ 是数据点数，$p$ 是拟合参数个数。

教学上最重要的直觉是：

\begin{itemize}
    \item 若 $\chi^2_\nu \gg 1$，模型可能没有解释主要结构，或误差被低估；
    \item 若 $\chi^2_\nu \ll 1$，误差可能被高估，或模型在当前数据上过于灵活；
    \item 若 $\chi^2_\nu$ 接近 $1$，说明模型与误差量级大致一致，但仍不代表模型必然物理正确。
\end{itemize}

\section{参数搜索：从网格到最佳模型}

在入门阶段，我们不必一开始就依赖复杂优化器。对只有一个主要参数 $\tau$ 的指数衰减模型，可以先用网格搜索建立最小算法直觉：

\begin{enumerate}
    \item 给定一组候选 $\tau$；
    \item 对每个 $\tau$ 计算模型曲线 $V(t;\tau)$；
    \item 计算对应的加权 $\chi^2(\tau)$；
    \item 选择 $\chi^2$ 最小的候选作为当前最佳估计；
    \item 再检查残差是否支持这个模型解释。
\end{enumerate}

如果模型还包含 $V_0$，可以把它一起放入网格，也可以在固定 $\tau$ 时求最优幅度。无论采用哪种实现，教学上最重要的是：\textbf{参数估计不是魔法函数给出的数字，而是模型预测和观测数据反复比较后的结果。}

\paragraph{固定 \texorpdfstring{$\tau$}{tau} 时怎样求最优 \texorpdfstring{$V_0$}{V0}}

指数模型可以写成

\[
V_i \approx V_0 a_i,
\qquad
a_i = e^{-t_i/\tau}.
\]

如果 $\tau$ 已经固定，未知量只剩下 $V_0$。这时加权卡方为

\[
\chi^2(V_0\mid\tau)
= \sum_i \frac{(V_i - V_0 a_i)^2}{\sigma_i^2}.
\]

对 $V_0$ 求导并令导数为零，可以得到

\[
\hat{V}_0(\tau)
=
\frac{\sum_i V_i a_i/\sigma_i^2}
{\sum_i a_i^2/\sigma_i^2}.
\]

这一步很重要，因为它把二维搜索压缩成了一维搜索：外层扫描 $\tau$，内层直接算出该 $\tau$ 下最优的 $V_0$。许多更复杂的拟合算法也在做类似事情：把一部分参数通过解析或线性代数求出，把真正困难的参数留给外层优化。

\begin{examplebox}[固定 tau 时估计 V0]
\begin{lstlisting}[style=pythonstyle]
def best_v0_for_tau(times, voltages, sigmas, tau):
    basis = [math.exp(-time / tau) for time in times]
    numerator = sum(v * a / (s ** 2) for v, a, s in zip(voltages, basis, sigmas))
    denominator = sum((a ** 2) / (s ** 2) for a, s in zip(basis, sigmas))
    return numerator / denominator
\end{lstlisting}
\end{examplebox}

\begin{examplebox}[一维参数搜索的最小形状]
\begin{lstlisting}[style=pythonstyle]
best = None
for tau in candidate_taus:
    v0 = best_v0_for_tau(times, voltages, sigmas, tau)
    predictions = [v0 * math.exp(-time / tau) for time in times]
    chi2 = sum(((obs - pred) / sigma) ** 2
               for obs, pred, sigma in zip(voltages, predictions, sigmas))
    if best is None or chi2 < best["chi2"]:
        best = {"tau": tau, "v0": v0, "chi2": chi2}
\end{lstlisting}
\end{examplebox}

\section{残差图应该怎样读}

残差分析的第一步不是计算更多指标，而是看残差是否还有结构。对指数衰减实验来说，可以特别检查：

\begin{itemize}
    \item 早期残差是否整体同号，暗示初始条件或响应延迟有问题；
    \item 晚期残差是否随时间系统偏离，暗示背景偏置或噪声模型不对；
    \item 残差大小是否明显超过误差条，暗示 $\sigma_i$ 可能被低估；
    \item 残差是否呈随机上下波动，说明当前模型至少解释了主要趋势。
\end{itemize}

\begin{warnbox}[好看的曲线不等于好模型]
一条曲线看起来穿过数据点，只能说明视觉上相近；残差是否有结构、参数是否有单位、模型是否来自机制关系，才决定它能否作为物理解释。
\end{warnbox}

\section{为什么模拟是实验分析的一部分}

许多初学者把“实验数据分析”和“数值模拟”当成两件完全分开的事，但在科学实践中，两者往往是一条闭环：

\begin{enumerate}
    \item 用理论或机制写出模型；
    \item 用模型前向生成预测曲线；
    \item 与实验数据比较；
    \item 根据偏差修正参数、假设或实验设置。
\end{enumerate}

对 RC 放电而言，最简单的前向模拟甚至可以从离散更新开始：

\[
V_{n+1} = V_n - \frac{\Delta t}{\tau} V_n.
\]

当 $\Delta t$ 足够小时，这个离散模型会逼近连续的指数衰减解。教学上，这能帮助学生理解“解析解”和“数值推进”之间的联系。

\begin{examplebox}[最小前向模拟]
\begin{lstlisting}[style=pythonstyle]
def forward_euler_decay(v0, tau, dt, n_steps):
    values = [v0]
    for _ in range(n_steps):
        values.append(values[-1] - (dt / tau) * values[-1])
    return values
\end{lstlisting}
\end{examplebox}

\section{为什么不能盲目做对数线性化}

面对指数衰减，很多学生会自然想到对公式两边取对数：

\[
\ln V = \ln V_0 - \frac{t}{\tau}.
\]

这样看起来就把非线性问题变成了线性拟合问题。这个思路在某些场景下当然有用，但教学上必须提醒两点：

\begin{itemize}
    \item 对数变换会改变误差结构，原本在 $V$ 空间近似对称的误差，在 $\ln V$ 空间里未必仍然对称；
    \item 当晚期电压很低、噪声相对占比变大时，取对数会显著放大不稳定性。
\end{itemize}

因此，\textbf{“线性化之后更方便”并不自动等于“统计上更合理”。} 这也是为什么本章 notebook 仍然优先在原始电压空间里做模型比较。

\section{Monte Carlo：为什么单次拟合还不够}

即使我们得到了一组最佳拟合参数，也仍然要追问：如果观测值在误差范围内稍微波动，参数会不会显著变化？这就是 Monte Carlo 重采样的动机。

最常见的教学做法是：

\begin{enumerate}
    \item 以原始测量值为中心，以 $\sigma_i$ 为尺度产生扰动样本；
    \item 对每个扰动样本重新拟合参数；
    \item 统计参数分布的均值、标准差和分位区间。
\end{enumerate}

这样得到的不是单个数字，而是一组带不确定度的参数估计。

\subsection{Monte Carlo 估计的是哪类不确定度}

这里也要保持概念清晰。当前章节中的 Monte Carlo 主要是在回答：

\begin{quote}
如果观测值在当前误差条附近发生随机扰动，拟合参数会波动多大？
\end{quote}

它主要覆盖的是\textbf{统计波动}，而不自动覆盖：

\begin{itemize}
    \item 模型形式本身错了怎么办；
    \item 背景系统误差或仪器漂移怎么办；
    \item 数据根本没有覆盖某种新行为怎么办。
\end{itemize}

因此，Monte Carlo 很重要，但它仍然只是可信度分析中的一层，不是全部答案。

\subsection{最小 Monte Carlo 代码骨架}

Monte Carlo 的核心是“扰动数据，重复完整拟合流程”，而不是只在最佳参数附近随便加一个误差条。最小代码骨架可以写成：

\begin{examplebox}[扰动数据并重复拟合]
\begin{lstlisting}[style=pythonstyle]
import random

tau_samples = []
for _ in range(n_trials):
    perturbed = [
        random.gauss(value, sigma)
        for value, sigma in zip(voltages, sigmas)
    ]
    best_tau = fit_tau_by_grid(times, perturbed, sigmas, candidate_taus)
    tau_samples.append(best_tau)

tau_mean = sum(tau_samples) / len(tau_samples)
\end{lstlisting}
\end{examplebox}

这段伪代码故意把 \texttt{fit\_tau\_by\_grid} 保留为一个函数名，因为教学重点是流程：每一次扰动后都要重新走完整参数估计，而不是只对最后结果做表面修饰。

\section{一个最小实验拟合结果报告模板}

一个更像科研报告的拟合结果，不应该只写“最佳 $\tau=2.15\,\mathrm{s}$”。至少应包含：

\begin{itemize}
    \item 最佳参数及其单位；
    \item 参数不确定度或 Monte Carlo 分布摘要；
    \item 使用的数据范围、误差模型和自由度；
    \item 残差是否存在系统结构；
    \item 与简单 baseline 或经验模型的比较。
\end{itemize}

例如，教学报告可以写成：“在原始电压空间中，用指数衰减模型拟合得到 $\tau \approx 2.15\,\mathrm{s}$；Monte Carlo 重采样给出的 $\tau$ 标准差约为若干百分之一秒；残差未显示强系统趋势，但该结论仍依赖当前误差条和理想 RC 模型假设。”这种写法比孤立参数更可信，因为它把数值、误差、诊断和适用边界放在了一起。到了后续项目证据包，原始实验字段和误差模型进入 Data Card，拟合流程和随机扰动方法进入 Evidence Record，残差与 baseline 比较进入 Model Experiment Record，物理解释和模型假设则进入 Trust Statement。

\section{本章 Evidence Record 附加字段}

本章不另建独立模型拟合记录模板，而是使用统一 Evidence Record，并补充模型、参数和残差字段：

\begin{examplebox}[Ch15 Evidence Record 附加字段]
\begin{lstlisting}[style=terminal]
Ch15 Fields:
- model formula:
- parameter meanings:
- parameter units:
- fitting method:
- residual plot:
- error metric / chi-square:
- uncertainty estimate:
- baseline comparison:
- physical interpretation limit:
\end{lstlisting}
\end{examplebox}

本章的特殊通过条件是：必须包含残差或误差诊断；不能只给拟合曲线。

\section{为什么还要和数据驱动近似比较}

如果你只追求插值效果，那么一个二次多项式或更灵活的回归器也可能在局部拟合得不错。但物理模型与数据驱动模型的差别在于：

\begin{itemize}
    \item 物理模型参数可解释，例如 $\tau$ 有明确单位和机制含义；
    \item 物理模型更容易外推到未观测区间并保持机制一致性；
    \item 数据驱动近似在局部可能更灵活，但未必尊重系统的真实动力学。
\end{itemize}

因此，本章 notebook 会同时给出指数模型和简单二次多项式近似的比较，帮助学生建立“拟合质量”和“物理可解释性”并不完全等价的认识。

\section{从实验案例中应该学会什么工作顺序}

把本章案例压缩成最小流程，可以写成：

\begin{enumerate}
    \item 先写出候选物理模型与参数含义；
    \item 再检查数据单位、误差和时间范围；
    \item 在原始观测空间里做拟合；
    \item 报告残差和卡方，而不是只报一条曲线；
    \item 用 Monte Carlo 估计参数波动；
    \item 再把物理模型与更灵活的经验拟合作比较。
\end{enumerate}

这条顺序很重要，因为它避免了一个常见误区：还没搞清数据和模型关系，就急着换更复杂的拟合器。

\section{教学数据与 notebook 流程}

本章数据文件 \nolinkurl{physics_model_fit_demo.csv} 采用一个极小 RC 放电教学表。每一行包含：

\begin{itemize}
    \item \texttt{sample\_id}：实验序列标识；
    \item \texttt{time\_s}：时间；
    \item \texttt{voltage\_v}：测得电压；
    \item \texttt{sigma\_v}：测量误差；
    \item \texttt{split}：数据角色；
    \item \texttt{experiment\_type}：教学标签。
\end{itemize}

配套 notebook 会完成：

\begin{itemize}
    \item 读入实验数据并检查时间与电压范围；
    \item 在一组候选 $\tau$ 上搜索最佳指数衰减模型；
    \item 计算残差和加权 $\chi^2$；
    \item 用一个极简 Monte Carlo 重采样估计 $\tau$ 的不确定度；
    \item 构造一个二次多项式近似，和指数模型做对比；
    \item 若可用，再画出观测点、模型曲线和残差结构。
\end{itemize}

\section{常见错误}

\begin{warnbox}[实验拟合中的典型问题]
\begin{itemize}
    \item 只看拟合曲线是否顺眼，不检查残差；
    \item 忽略测量误差，默认每个点同样可信；
    \item 用高阶经验函数拟合后，误以为已经获得了物理解释；
    \item 报告最佳参数，却不给任何不确定度估计；
    \item 只会调用库函数，不知道模型公式来自哪条物理关系。
\end{itemize}
\end{warnbox}

\section{AI 助手如何帮助你}

在这一章，AI 助手适合帮助你：

\begin{itemize}
    \item 把模型拟合、残差和 Monte Carlo 代码组织得更清晰；
    \item 检查单位和参数定义是否一致；
    \item 帮你把指数衰减推导和数值模拟的联系讲得更明白；
    \item 协助你比较“物理模型拟合更好”与“经验模型拟合更灵活”之间的差异。
\end{itemize}

但某个模型假设是否真的适合实验系统、某个残差结构是否意味着新物理或系统误差，仍然需要你自己结合背景知识判断。

\section{练习}

\begin{exercisebox}[基础练习]
用本章数据计算每个观测点的残差，并说明残差在早期时间和晚期时间是否表现出系统差异。
\end{exercisebox}

\begin{exercisebox}[进阶练习]
从
\[
\frac{dV}{dt} = -\frac{1}{\tau}V
\]
出发，重新推导指数衰减解，并解释为什么 $\tau$ 可以看成系统衰减到原来约 $e^{-1}$ 倍的时间尺度。
\end{exercisebox}

\begin{exercisebox}[开放练习]
为一个你熟悉的物理实验写出“模型拟合前检查清单”，至少包含：变量单位、候选机制模型、误差来源、残差检查和不确定度估计方法。
\end{exercisebox}

\section{Exit Gate}

\begin{enumerate}
    \item 模型公式、参数和单位是否清楚？
    \item 拟合方法和误差指标是否说明？
    \item 是否查看残差或不确定度？
    \item 物理模型和 ML baseline 的关系是否说明？
\end{enumerate}

最低交付：Evidence Record、拟合结果和残差/不确定度检查。

\section{本章小结}

物理实验数据分析的关键，在于把观测点、机制模型和参数解释连成一条链。本章通过一个极小 RC 放电示例，展示了从微分方程、指数模型、卡方统计和 Monte Carlo 不确定度，到经验模型比较的完整流程。只要这条链条真正建立起来，你后面面对摆实验、衰变计数或更复杂数值模拟时，就不会只停留在“把曲线画出来”的层面。
